\documentclass[11pt]{article}
\usepackage[a4paper,margin=21mm]{geometry}
\usepackage[no-math]{fontspec}
\setmainfont{Latin Modern Roman}
\newfontfamily\CJKfont{Noto Serif CJK SC}
\XeTeXlinebreaklocale "zh"
\XeTeXlinebreakskip=0pt plus 1pt
\usepackage{amsmath,amssymb,amsthm,amscd,graphicx,color}
\usepackage{xurl}
\usepackage[unicode,hidelinks]{hyperref}
\allowdisplaybreaks
\setlength\emergencystretch{3em}
\numberwithin{equation}{section}

\numberwithin{equation}{section}

\newtheorem{Theorem}{Theorem}[section]
\newtheorem{Corollary}[Theorem]{Corollary}
\newtheorem{Lemma}[Theorem]{Lemma}
\newtheorem{Proposition}[Theorem]{Proposition}
\newtheorem{Conjecture}[Theorem]{Conjecture}
 { \theoremstyle{definition}
\newtheorem{Definition}[Theorem]{Definition}
\newtheorem{Example}[Theorem]{Example}
\newtheorem{Remark}[Theorem]{Remark} }

\newcommand{\R}{\mathbb{R}}
\newcommand{\C}{\mathbb{C}}
\newcommand{\Z}{\mathbb{Z}}
\newcommand{\Q}{\mathbb{Q}}
\newcommand{\g}{\mathfrak{g}}
\newcommand{\h}{\mathfrak{h}}


\begin{document}
\begin{center}\Large Finite-dimensional twisted traces and finite-parameter nondegenerate short star-products for conical symplectic quantizations with reductive graded Poisson automorphisms\end{center}
\noindent\textbf{Stable ID:} 1494\quad\textbf{Author:} Pavel Etingof; Douglas Stryker\par
\noindent\textbf{Work:} Short Star-Products for Filtered Quantizations, I\par
\noindent\textbf{Source:} \url{https://sigma-journal.com/2020/014/}\par
\noindent\textbf{Version:} Official SIGMA 16 (2020) article 014, DOI 10.3842/SIGMA.2020.014; journal PDF and native TeX acquired 2026-09-30, exact bytes pinned by SHA256\par
\noindent\textbf{Rights:} CC BY 4.0. Authors retain copyright. \url{https://creativecommons.org/licenses/by/4.0/}. Reuse and typesetting adaptation; original language and mathematical content preserved.\par
\noindent\textbf{Difficulty (prior selection preserved):} {\CJKfont H3: The proof identifies nondegenerate short products with filtered twisted traces and bounds their parameter space by fixed-scheme Poisson homology. The semisimple case uses an associated-graded commutator quotient and finite symplectic-leaf geometry. The nonsemisimple case requires a nilpotent Jordan degeneration and an injective leading-trace argument, rather than merely importing Poisson-homology finiteness.}\par
\medskip\noindent\textbf{Editorial boundary note (not author text):} Selected Theorem3.15 finite twisted-trace/short-star-product parameter bound for normal conical symplectic singularities with the stated reductive graded Poisson automorphism group, in the defined twisted-bimodule scope of an actual filtration-preserving automorphism g of the chosen parity-equivariant quantization. The original broader g wording remains literal; no blanket lift of arbitrary graded automorphisms or existence/positivity of short products is asserted. Full graded Poisson/filtered parity/quantization-map/short/nondegenerate conventions, inner-derivation and fixed-scheme arguments, complete twisted trace/orthogonal complement correspondence, semisimple associated-graded/Poisson-homology bound and entire nilpotent Jordan degeneration/leading-trace injection are retained. The two original Straightforward inverse-quantization-map comments have both explicit mutually inverse formulas retained; no generated routine expansion. Named finite-leaf Poisson homology, conical geometry, GAGA and Jacobson-Morozov results remain precise established references. Source symbols/formulas/order and all necessary footnotes retain original numbering. Hyperkahler/existence/positivity/examples and supporting units have no additional counts. All selected core is native editable TeX, with no assistant proof.\par
\medskip\hrule\medskip\noindent\textbf{Author text begins below.}\par
\setcounter{section}{1}
% Original source lines 89--94
The goal of this paper is to partially prove the conjectures from~\cite{BPR}.
Namely, let us say that a short star-product $*$ is {\it nondegenerate}
if the bilinear form $(a,b)={\rm CT}(a*b)$ on $A$ (where ${\rm CT}$ stands for the constant term)
 is nondegenerate. This is a natural condition from the physics point of view; e.g.,
it follows from the positivity condition, saying that the Hermitian version of this form is positive definite. Our first main result is that nondegenerate short star-products for hyperK\"ahler cones depend on finitely many parameters, confirming a conjecture from~\cite{BPR}.
The proof is based on the idea of Kontsevich to express nondegenerate short star-products in terms of {\it nondegenerate twisted traces}, which can then be understood as elements of a certain zeroth Hochschild homology group, as well as a number of results on symplectic singularities~\cite{Ka,L2,Na3}, hyperK\"ahler manifolds~\cite{HKLR} and Poisson homology~\cite{ES}.

\setcounter{footnote}{3}
% Original source lines 107--199
\section{Preliminaries}\label{section2}

\subsection{Star-product quantization of a graded Poisson algebra}
Let $A$ be a commutative ${\mathbb C}$-algebra with a $\Z_{\geq 0}$ grading given by
\[
A = \bigoplus_{d \geq 0} A_d.
\]
Let $\{~,~\}$ be a Poisson bracket of degree $-2$ on $A$; namely,
\[
\{~,~\} \colon \ A_n \otimes A_m \to A_{n+m-2}.
\]
A basic example is the symmetric algebra $A=SV$ of a finite dimensional symplectic vector space $V$
with the Poisson bracket defined by the symplectic form.

Note that this setting includes the case of graded Poisson algebras with Poisson bracket of degree $-1$ (such as the symmetric algebra of a Lie algebra, the algebra of regular functions on the nilpotent cone, etc.), by multiplying the degrees by $2$.\footnote{In fact, this theory can be developed for Poisson brackets of any negative degree, but for most applications it is sufficient to consider degree $-2$, so we restrict our attention to this case.}

\begin{Definition}
A \emph{star-product} on $A$ quantizing $\lbrace\,,\,\rbrace$ is an associative multiplication operation
\[
*\colon \ A \otimes A \to A
\]
such that for $a \in A_n$ and $b \in A_m$,
\[
a*b = \sum_{k = 0}^{\lfloor \frac{n+m}{2}\rfloor} C_k(a,b),
\]
where $C_k \colon A_n \otimes A_m \to A_{n+m-2k}$ are bilinear maps such that $C_0(a,b) = a b$ and
\begin{gather*}
C_1(a,b) - C_1(b,a) = \{a,b\}.
\end{gather*}
\end{Definition}

\begin{Definition} A star-product $*$ on $A$ is \emph{short} if
for any $m,n\in {\mathbb Z}_{\ge 0}$ and any $a \in A_n$, $b \in A_m$ one has $C_k(a,b) = 0$ for all $k > \min(n,m)$. In other words, $a * b$ has no terms in $A_d$ for $d < |n - m|$.
\end{Definition}

From now on assume that $A_0={\mathbb C}$ and $\dim A_i<\infty$ for all $i$. We define the following (in general, non-symmetric) inner product on $A$ induced by the star-product. Let $a \in A_n$ and $b \in A_m$. Then
\[
\langle a, b \rangle := \text{CT}(a * b),
\]
where $\text{CT}\colon A \to A_0$ is the constant term map, taking the term of degree zero. Note that if $*$ is short, then $\langle a, b \rangle = 0$ if $n \neq m$, i.e., the decomposition $A=\oplus_{d\ge 0}A_d$ is orthogonal.

\begin{Definition} A short star-product $*$ on $A$ is said to be {\it nondegenerate} if the corresponding inner product $\langle \,,\,\rangle$ is nondegenerate in each degree~$i$.
\end{Definition}

\subsection{Construction of star-products from a quantum algebra}\label{constru} Let $A$ be a ${\mathbb Z}_{\ge 0}$-graded algebra. Fix a filtered associative algebra
\[
{{\mathbf A}} = \bigcup_{d \geq 0} F_d{{\mathbf A}}
\]
such that its associated graded algebra is identified with the graded algebra $A$; namely,
\[
\operatorname{gr}{{\mathbf A}} = \bigoplus_{d \geq 0} F_{d}{{\mathbf A}}/F_{d-1}{{\mathbf A}}
\]
and $F_{d}{{\mathbf A}}/F_{d-1}{{\mathbf A}} \cong A_d$ for $d \geq 0$ (compatibly with multiplication) with $F_{-1}{{\mathbf A}} := 0$.
Assume that we have a filtration preserving ${\mathbb Z}/2$-action $s\colon {{\mathbf A}}\to {{\mathbf A}}$ such that the corresponding associated graded map $s\colon A\to A$ is given by $s=(-1)^d$, where $d\colon A\to A$ is the degree operator. In this case, it is easy to see that $[F_n{{\mathbf A}},F_m{{\mathbf A}}]\subset F_{n+m-2}{{\mathbf A}}$, so we have the leading coefficient map $\lbrace\,,\,\rbrace\colon A_n\otimes A_m\to A_{n+m-2}$. It is easy to check that $\lbrace\,,\,\rbrace$ is a Poisson bracket on $A$, and ${{\mathbf A}}$ is a~${\mathbb Z}/2$-equivariant filtered quantization of $\lbrace\,,\,\rbrace$.\footnote{Since all quantizations we will consider will be ${\mathbb Z}/2$-equivariant, we will not mention it explicitly from now on.}

\begin{Definition}A \emph{quantization map} is a ${\mathbb Z}/2$-equivariant isomorphism of vector spaces
\[
\phi \colon \ A \to {{\mathbf A}}
\]
such that $\phi(A_d) \subseteq F_d{{\mathbf A}}$ for $d\ge 0$ and $\operatorname{gr} \phi \colon A \to A$ is the identity map, where
\[
\operatorname{gr} \phi = \bigoplus_{d\geq 0} \operatorname{gr} \phi_d, \qquad \operatorname{gr} \phi_d \colon \ A_d \to F_d{{\mathbf A}}/F_{d-1}{{\mathbf A}}=A_d.
\]
\end{Definition}
Note that $\phi$ need not (and usually cannot) be a homomorphism of algebras, since the algebra~$A$ is commutative but ${{\mathbf A}}$ is not in general.

Given a quantization map $\phi\colon A \to {{\mathbf A}}$, we can define the star-product on $A$ by
\[
a * b := \phi^{-1}(\phi(a)\phi(b)).
\]
Conversely, given a star-product on $A$, we can set ${{\mathbf A}}:=(A,*)$ with the filtration induced by the grading, and $\phi:={\rm Id}$.

\begin{Proposition} This defines a pair of mutually inverse bijections between star-products on~$A$ and ${\mathbb Z}/2$-equivariant quantizations of $A$ equipped with a quantization map.
\end{Proposition}

\begin{proof} Straightforward.
\end{proof}

\subsection{Even star-products}

\begin{Definition} A star-product $*$ on $A$ is said to be {\it even} if $C_k(a,b)=(-1)^k C_k(b,a)$ for all $k\ge 0$.
\end{Definition}

In particular, for even star-products $C_1$ is skew-symmetric, so since $C_1(a,b)-C_1(b,a)=\lbrace a,b\rbrace$, we have $C_1(a,b)=\frac{1}{2}\lbrace a,b\rbrace$.

Let $*$ be an even star-product on $A$ and ${{\mathbf A}}=(A,*)$ be the corresponding quantum algebra.
Then the map $\sigma={\rm i}^d$ defines an antiautomorphism ${{\mathbf A}}\to {{\mathbf A}}$ such that $\sigma^2=s=(-1)^d$. A ${\mathbb Z}/2$-invariant quantization ${\mathbf A}$ of $A$ equipped with a~filtration-preserving antiautomorphism $\sigma$ such that $\sigma^2=s$ will be called {\it even}. Thus an even star-product gives rise to an even quantization. Conversely, an even quantization ${{\mathbf A}}$ of $A$ gives rise to an even star-product on~$A$.

\begin{Proposition} This defines a pair of mutually inverse bijections between even star-products on $A$ and ${\mathbb Z}/2$-equivariant quantizations ${{\mathbf A}}$ of $A$ equipped with a filtered antiautomorphism $\sigma\colon {{\mathbf A}}\to {{\mathbf A}}$ such that $\operatorname{gr}\sigma={\rm i}^d$ and $\sigma^2=s$ with a $\sigma$-invariant $($i.e., ${\mathbb Z}/4$-invariant$)$ quantization map.
\end{Proposition}

\begin{proof} Straightforward.
\end{proof}

\setcounter{Theorem}{8}
% Original source lines 231--254
\subsection{Conical symplectic singularities}

The main class of examples of algebras $A$ we will be interested in is $A={\mathcal O}(X)$, where~$X$ is a (normal) conical symplectic singularity in the sense of Beauville (see~\cite{Ka, L2} and references therein for their definition and basic properties).

We will need the following lemma.

\begin{Lemma}[{\cite[Proposition~2.5]{L2}}]\label{inn} If $X$ is a conical symplectic singularity then every Poisson derivation of $A:=\mathcal{O}(X)$ is inner.
\end{Lemma}

\begin{proof} Here is another proof. Any such homogeneous derivation is defined by an (a priori multivalued) holomorphic Hamiltonian $H$ on the big symplectic leaf $X^\circ$ of $X$. Moreover, $H$ is, in fact, single-valued, since by~\cite{Na1} the algebraic fundamental group of $X^\circ$ is finite (indeed, we have a homomorphism $\phi\colon \pi_1(X^\circ,x_0)\to {\mathbb C}$ given by $\phi(\gamma)=\gamma(H)-H$, where $\gamma(H)$ is the analytic continuation of~$H$ along~$\gamma$, and $\phi=0$ since $\pi_1(X^\circ,x_0)_{\rm alg}$ is finite). Then by normality of~$X$ the function~$H$ extends to the whole~$X$, to a homogeneous holomorphic function on~$X$. This function is a holomorphic section of a line bundle on ${\mathbb P}X$, which is algebraic by the GAGA theorem, so $H\in \mathcal{O}(X)$, as claimed.
\end{proof}

It is shown in \cite{L2} that for conical symplectic singularities there is a very nice parametrization of filtered quantizations of $A$, which is the same as Namikawa's parametrization of filtered Poisson deformations~$X$~\cite{Na3}. Namely, they are parametrized by the finite dimensional space $\mathfrak{P}=H^2\big(\widetilde X^{\rm reg},{\mathbb C}\big)$, where $\widetilde X^{\rm reg}$ is the smooth part of the ${\mathbb Q}$-terminalization $\widetilde X$ of~$X$, modulo the action of a finite real reflection group~$W$ called the Namikawa Weyl group. Denote the quantization corresponding to $\lambda\in \mathfrak{P}$ by~${\mathbf A}_\lambda$. It is easy to show that ${\mathbf A}_\lambda$ is ${\mathbb Z}/2$-invariant.

Since $W$ is a reflection group, the space $\mathfrak{P}/W$ bijectively parametrizing quantizations is an affine space, whose coordinates are independent homogeneous generating invariants $p_1,\dots,p_r\in {\mathbb C}[\mathfrak{P}]^W$ of degrees $d_1,\dots,d_r$.

It is not hard to see that ${\mathbf A}_\lambda$ is even if and only if $\lambda$ and $-\lambda$ are equivalent with respect to the Namikawa Weyl group, i.e., if and only if there is an element $w\in W$ such that $w\lambda=-\lambda$ (see \cite{BPR}). This is equivalent to saying that $p_i=0$ whenever $d_i$ is odd, which defines a subspace in
$\mathfrak{P}/W$.

\begin{Remark}Lemma \ref{inn} implies that in the case of conical symplectic singularities the antiautomorphism $\sigma\colon {\mathbf A}\to {\mathbf A}$ such that $\sigma^2=s$ and $\operatorname{gr}\sigma={\rm i}^d$ is unique if exists. Indeed, if~$\sigma_1$,~$\sigma_2$ are two such antiautomorphisms then $g:=\sigma_1\circ \sigma_2^{-1}$ is an automorphism commuting with~$s$ such that $\operatorname{gr} g=1$. Hence $\log (g)$ is a linear combination of derivations of~${\mathbf A}$ of degree $\le -2$. If $g\ne 1$, then in the leading order $\log(g)$ yields a nonzero homogeneous derivation $D\colon A\to A$ of degree $\le -2$. By Lemma~\ref{inn}, it is inner, so given by a Hamiltonian of degree $\le 0$. But such a~Hamiltonian is necessarily constant, so $D=0$, a contradiction.

The same argument proves that if $g\colon {\mathbf A}\to {\mathbf A}$ is a filtration-preserving
automorphism commuting with $s$ then $g$ is completely determined by $\operatorname{gr}g$.
\end{Remark}

\setcounter{subsection}{5}\setcounter{footnote}{6}
% Original source lines 268--351
\subsection{Finiteness of the number of fixed sets}

\begin{Lemma}\label{finfix} Let $X$ be an affine scheme of finite type over ${\mathbb C}$ and
$G$ be a reductive group acting on $X$. Then up to the $G$-action there are finitely many subschemes
of $X$ of the form $X^g$, where $g\in G$ is a semisimple element, .
\end{Lemma}

\begin{proof} First consider the case when $G$ is diagonalizable (in which case every $g\in G$ is semisimple). In this case the lemma is standard (see, e.g., \cite[Proposition on p.~44]{Ste}). Namely, pick generators $f_1,\dots,f_n$ of ${\mathbb C}[X]$ such that $f_i(gx)=\chi_i(g)f_i(x)$ for some characters~$\chi_i$ of~$G$. Then $f_i(gx)-f_i(x)=(\chi_i(g)-1)f_i(x)$, so~$X^g$ is cut out by the equations $f_i(x)=0$ for those~$i$ for which $\chi_i(g)\ne 1$, which implies that there are finitely many subschemes~$X^g$.

The general case reduces to the diagonalizable case by using the following fact (see~\cite{Sie} and \cite[Section~1.14]{Lu}): for each connected component~$Z$ of~$G$ there is a diagonalizable subgroup $D_Z\subset G$ such that every semisimple element $g\in Z$ is $G^0$-conjugate to an element of~$D_Z$.\footnote{We are grateful to G.~Lusztig for this explanation.}
\end{proof}

\section{Classification of short star-products}\label{section3}

In this section we present the classification of nondegenerate short star-products, the idea of which was explained to the first author by Maxim Kontsevich.

\subsection{Twisted traces}

Consider the setting of Section~\ref{constru}. Let $g\colon {{\mathbf A}}\to {{\mathbf A}}$ be a filtration-preserving invertible linear map, and let $T\colon {{\mathbf A}}\to {\mathbb C}$ be a $g$-twisted trace, i.e., $T({\mathbf a}{\mathbf b})=T({\mathbf b}g({\mathbf a}))$ for all ${\mathbf a},{\mathbf b}\in {{\mathbf A}}$. Note that taking ${\mathbf b}=1$, we get $T(g({\mathbf a}))=T({\mathbf a})$, i.e., $T$ is $g$-invariant.

Define the bilinear form $(\,,\,)_T$ on ${\mathbf A}$ by the formula $({\mathbf a},{\mathbf b})_T:=T({\mathbf a}{\mathbf b})$.

\begin{Lemma}\label{ide} Let $K$ be the left kernel of the bilinear form $(\,,\,)_T$. Then $K$ is also the right kernel of $(\,,\,)_T$, and a $g$-invariant two-sided ideal in~${{\mathbf A}}$. In particular, if ${\mathbf A}$ is a simple algebra and $T\ne 0$ then $K=0$.
\end{Lemma}

\begin{proof} Let $K'$ be the right kernel of $(\,,\,)_T$. We have ${\mathbf b}\in K'$ if and only if $T({\textbf{ab}})=0$ for all ${\mathbf a}\in {{\mathbf A}}$. This is equivalent to saying that $T({\mathbf b}g({\mathbf a}))=0$ for all ${\mathbf a}\in {{\mathbf A}}$, which happens if and only if ${\mathbf b}\in K$. Thus, $K'=K$. Also ${\mathbf a}\in K$ if and only if $T({\mathbf b}g({\mathbf a}))=0$ for all ${\mathbf b}\in {\mathbf A}$, which happens if and only if $g({\mathbf a})\in K$, hence $K$ is $g$-invariant. Finally, if ${\mathbf a}\in K$ and ${\mathbf b}\in {{\mathbf A}}$
then $({\textbf{ab}},{\textbf{c}})_T= ({\textbf{a}},{\textbf{bc}})_T=0$ and $({\textbf{c}},{\textbf{ba}})_T=({\textbf{cb}},{\textbf{a}})_T=0$, so ${\textbf{ab}},{\textbf{ba}}\in K$, i.e., $K$ is a two-sided ideal, as desired.
\end{proof}

By Lemma~\ref{ide}, it makes sense to speak of the kernel of $(\,,\,)_T$ (without specifying if it is left or right).

\begin{Lemma} If the kernel of $(\,,\,)_T$ vanishes then $g$ is an algebra automorphism.
\end{Lemma}

\begin{proof}We have
\begin{gather*}
T({\mathbf c}g({\mathbf a}{\mathbf b}))=T({\textbf{abc}})=T({\textbf{bc}}g({\mathbf a}))=T({\mathbf c}g({\mathbf a})g({\mathbf b})),
\end{gather*}
hence $g({\mathbf{ab}})=g({\mathbf a})g({\mathbf b})$.
\end{proof}

\subsection{Nondegenerate short star-products and nondegenerate twisted traces}

We will say that $T$ is {\it nondegenerate} if the restriction of the bilinear form $({\mathbf a},{\mathbf b})_T:=T({\mathbf a}{\mathbf b})$ to~$F_i{{\mathbf A}}$ is nondegenerate for each~$i$. Note that this is a strictly stronger condition than vanishing of the kernel of $(\,,\,)_T$ (for example, it implies that $T(1)\ne 0$).

Suppose that $T$ is nondegenerate and $s$-invariant (in this case $g$ commutes with~$s$). Then the bilinear form $(\,,\,)_T$ defines a canonical orthogonal complement ${\mathbf A}^T_i$ to $F_{i-1}{{\mathbf A}}$ in $F_i{{\mathbf A}}$. Note that the right and left orthogonal complement coincide since $g$ is filtration-preserving. We have a natural isomorphism $\theta_i^T\colon {\mathbf A}_i^T\to A_i$. Define the quantization map $\phi^T\colon A\to {{\mathbf A}}$ by $\phi^T|_{A_i}:=\big(\theta_i^T\big)^{-1}$. Clearly, $\phi^T$ is $s$-invariant. Thus, $\phi^T$ defines a star-product $*$ on $A$. Clearly, this star-product does not depend on the normalization of~$T$, so we will always normalize~$T$ so that $T(1)=1$ (which is possible since for a nondegenerate~$T$, one has $T(1)\ne 0$).

\begin{Proposition}[M.~Kontsevich]\label{kon} The star-product $*$ is short and nondegenerate. Moreover, any nondegenerate short star-product on~$A$ is obtained in this way from a unique nondegenerate~$T$ such that $T(1)=1$. In other words, this correspondence is a bijection between short nondegenerate star-products on $A$ and filtered quantizations ${{\mathbf A}}$ of~$A$ equipped with a filtration-preserving automorphism~$g$ commuting with $s$ and a nondegenerate $s$-invariant $g$-twisted trace~$T$ such that $T(1)=1$.
\end{Proposition}

\begin{proof} We have a nondegenerate (in general, nonsymmetric) inner product on~$A$ given by $\beta(a,b):=(\phi(a),\phi(b))_T$, and the spaces $A_i$ are orthogonal under this inner product. Moreover, it is clear that $\beta(a*b,c)=\beta(a,b*c)$. Suppose $a\in A_m$, $b\in A_n$. Let us show that every homogeneous component of $a*b$ has degree $\ge |m-n|$. First assume that $m>n$. It suffices to show that for any homogeneous $c\in A$ of degree $d<m-n$, one has $\beta(a*b,c)=0$. But we have $\beta(a*b,c)=\beta(a,b*c)$, and $b*c$ does not have any terms of degree bigger than $n+d<m$.
Hence $\beta(a*b,c)=\beta(a,b*c)=0$, as desired. A similar argument applies if $m<n$: we have $\beta(c,a*b)=\beta(c*a,b)=0$. Thus, the star-product~$*$ is short. Moreover, we have
\begin{gather*}
\langle a,b\rangle={\rm CT}(a*b)=\beta(a*b,1)=\beta(a,b),
\end{gather*}
which implies that $\langle\,,\,\rangle$ (and hence~$*$) is nondegenerate.

Conversely, assume that $*$ is a nondegenerate short star-product on~$A$. Then the quantization~${{\mathbf A}}$ is identified with $(A,*)$ using the quantization map~$\phi$. Set $T({\mathbf a}):={\rm CT}\big(\phi^{-1}({\mathbf a})\big)$. Then the form $({\textbf{a}},{\textbf{b}})_T:=T({\textbf{ab}})$ is nondegenerate on each $F_i{{\mathbf A}}$. Thus there exists a unique linear automorphism~$g_i$ of $F_i{{\mathbf A}}$ such that $T({\textbf{ab}})=T({\mathbf b}g_i({\mathbf a}))$. Moreover, it is clear that~$g_i$ preserves~$F_{i-1}{{\mathbf A}}$, and $g_i|_{F_{i-1}{{\mathbf A}}}=g_{i-1}$. Thus, the maps~$g_i$ define a~filtration-preserving linear automorphism $g\colon {{\mathbf A}}\to {{\mathbf A}}$.

Finally, it is easy to check that these assignments are mutually inverse, as claimed.
\end{proof}

\begin{Remark} If $T$ is any $s$-invariant linear functional on ${\mathbf A}$ such that the form $({\mathbf a},{\mathbf b})_T=T({\mathbf a}{\mathbf b})$ is nondegenerate on each $F_i{\mathbf A}$ (which clearly happens for ``random'' $T$) then for each~$i$ we can define a linear automorphism $g_i$ of $F_i({\mathbf A})$ such that $T({\textbf{ab}})=T({\mathbf b}g_i({\mathbf a}))$. However, in general $g_i|_{F_{i-1}{\mathbf A}}\ne g_{i-1}$ and $g_i$ does not preserve~$F_{i-1}{\mathbf A}$ (so the automorphism~$g$ of ${\mathbf A}$ restricting to~$g_i$ on each $F_i{\mathbf A}$ is not defined). Thus, the left and right orthogonal complements of~$F_{i-1}{\mathbf A}$ in~$F_i{\mathbf A}$ do not coincide, and one cannot define a quantization map giving a short star-product. The condition that~$g$ is well defined is a very strong restriction of~$T$, and we will see that it usually leads to a finite dimensional space of possible~$T$.
\end{Remark}

Let ${\mathbf A}g$ be the ${\mathbf A}$-bimodule which is ${\mathbf A}$ as a left module and ${\mathbf A}$ with action twisted by $g$ as a right module.

\begin{Corollary}\label{bun} Let $S=S({\mathbf A})$ be the set of nondegenerate short star-products corresponding to a filtered quantization ${{\mathbf A}}$ of $A$. Then we have a map $\pi\colon S\to \operatorname{Aut}({{\mathbf A}})$ from $S$ to the group of filtration-preserving $s$-invariant automorphisms of ${{\mathbf A}}$ whose fiber $\pi^{-1}(g)$ is a subset of the space $(HH_0({{\mathbf A}},{{\mathbf A}}g)^s)^*$ dual to the $s$-invariants in the Hochschild homology $HH_0({{\mathbf A}},{{\mathbf A}}g)$.
\end{Corollary}

\begin{proof} This follows immediately from Proposition~\ref{kon}, since by definition, an $s$-invariant $g$-twisted trace on ${{\mathbf A}}$ is an element of $(HH_0({{\mathbf A}},{{\mathbf A}}g)^s)^*$.
\end{proof}

\begin{Corollary} \label{bun1} Suppose $A=\mathcal{O}(X)$, where $X$ is a conical symplectic
singularity. Then any nondegenerate short star-product corresponding to $g\in \operatorname{Aut}({{\mathbf A}})$
is invariant with respect to the connected component~$Z_g^0$ of the centralizer~$Z_g$
of~$g$ in $\operatorname{Aut}({{\mathbf A}})$. In particular, any nondegenerate even short star-product
is invariant under $\operatorname{Aut}({\mathbf A})^0$.
\end{Corollary}

\begin{proof} It suffices to show that the Lie algebra $\operatorname{Lie}Z_g$ acts trivially on
$HH_0({\mathbf A},{\mathbf A} g)$. Let $z\in \operatorname{Lie}Z_g$. By Lemma~\ref{inn}, there exists a unique ${\mathbf z}\in F_2{\mathbf A}$ up to adding a constant such that $z({\mathbf a})=[{\mathbf z},{\mathbf a}]$ for all ${\mathbf a}\in {\mathbf A}$. Then $g({\mathbf z})={\mathbf z}+C$ for some constant $C$, so $C=g({\mathbf z})-{\mathbf z}$. Thus if $C\ne 0$ then for any $g$-twisted trace $T$ we have $T(1)=0$, so there are no nondegenerate short star-products at all. On the other hand, if $C=0$ then we have $[{\mathbf z},{\mathbf a}]={\mathbf z}{\mathbf a}-{\mathbf a}g({\mathbf z})$, so ${\mathbf z}$ acts trivially on $HH_0({\mathbf A},{\mathbf A} g)$, as claimed. The last statement follows from the fact that in the even case $g=s$, and~$s$ by definition is central in~$\operatorname{Aut}({{\mathbf A}})$.
\end{proof}


\setcounter{subsection}{4}\setcounter{Theorem}{10}\setcounter{footnote}{8}
% Original source lines 460--491
\subsection{Finite dimensionality of the space of nondegenerate short star-products}

\subsubsection{The case of a semisimple automorphism}

\begin{Proposition}\label{finde1} Let $A$ be a finitely generated Poisson algebra such that the Poisson scheme $X=\operatorname{Spec}A$ has finitely many symplectic leaves. Let~${\mathbf A}$ be a quantization of~$A$, and \mbox{$g\colon {{\mathbf A}}\to {{\mathbf A}}$} be a semisimple automorphism which commutes with~$s$. Then nondegenerate short star-products on~$A$, if they exist, are parametrized by points of the finite dimensional vector space \linebreak $(HH_0({{\mathbf A}},{{\mathbf A}}g)^s)^*$ which don't belong to a countable collection of hypersurfaces, modulo scaling.

Moreover, $\dim HH_0({{\mathbf A}},{{\mathbf A}}g)^s\le \dim HP_0(X^g)^s<\infty$, where $X^g$ is the fixed point scheme of~$g$ on~$X$, and $HP_0(X^g):= \mathcal{O}(X^g)/\lbrace{\mathcal{O}(X^g), \mathcal{O}(X^g)\rbrace}$ is the zeroth Poisson homology of $X^g$.\footnote{Here, abusing notation, we denote the associated graded isomorphism $\operatorname{gr}g\colon A\to A$ simply by $g$.}
\end{Proposition}

\begin{proof} Let $E$ be the span of elements ${\mathbf{ab}}-{\mathbf{b}}g({\mathbf a})$ for ${\mathbf{a}},{\mathbf{b}}\in {{\mathbf A}}$.

We claim that for any $a,b\in A$, $\operatorname{gr}(E)$ contains the element $(a-g(a))b$. Indeed, we may assume that $a$, $b$ are homogeneous of degrees $i$, $j$, and let ${{\mathbf a}}$, ${\mathbf b}$ be their lifts to ${{\mathbf A}}$. Then the leading term (of degree $i+j$) of ${{\mathbf a}} {\mathbf b}-{\mathbf b}g({{\mathbf a}})$ is $(a-g(a))b$, as desired.

\looseness=-1 Also, we claim that if $a\in A$ and $g(a)=a$ then for any $b\in A$, $\lbrace{a,b\rbrace}$ belongs to $\operatorname{gr}(E)$. Indeed, we can again assume that $a$, $b$ are homogeneous of degrees $i$, $j$. Since $g=1$ on its generalized 1-eigenspace, we can choose a lift ${\mathbf a}$ of~$a$ to~${\mathbf A}$ which is invariant under~$g$. Then choosing any lift~${\mathbf b}$ of~$b$ to~${\mathbf A}$, we see that the leading term (of degree $i+j-2$) of ${{\mathbf a}} {\mathbf b}-{\mathbf b}g({{\mathbf a}})=[{{\mathbf a}},{\mathbf b}]$ is $\lbrace{ a,b\rbrace}$, as desired.

Now let $I\subset A$ be the span of elements $(a-g(a))b$. This is the ideal of the fixed point subscheme~$X^g$. Note that $I\cap A^g$ is a Poisson ideal in $A^g$. Hence, $A/I=A^g/(I\cap A^g)$ is a Poisson algebra, so~$X^g$ is a Poisson scheme. Moreover, since $g=1$ on its generalized 1-eigenspace, the $g$-invariants in the tangent space to the symplectic leaf~$L$ at any point $P\in X^g\cap L$ form a~nondegenerate subspace, hence $X^g\cap L$ is symplectic. Thus, $X^g$ also has finitely many symplectic leaves, which are connected components of intersections of $X^g\cap L$ for various~$L$.

 Now, the space $A/\operatorname{gr}(E)$ receives a surjective $s$-invariant map from the space $HP_0(X^g)=\mathcal{O}(X^g)/\lbrace{ \mathcal{O}(X^g),\mathcal{O}(X^g) \rbrace}$. This space is finite dimensional by \cite[Theorem~1.4]{ES}. Hence the space $A/\operatorname{gr}(E)$
 is finite dimensional. Therefore the space ${{\mathbf A}}/E=HH_0({{\mathbf A}},{{\mathbf A}}g)$ is finite dimensional, and we have the inequalities
\begin{gather*}
\dim HH_0({{\mathbf A}},{{\mathbf A}}g)\le \dim HP_0(X^g),\qquad \dim HH_0({{\mathbf A}},{{\mathbf A}}g)^s\le \dim HP_0(X^g)^s.
\end{gather*}

\looseness=-2 Now the proposition follows, since by Corollary~\ref{bun}, short nondegenerate star-products corre\-spon\-ding to~$g$ are parametrized by the subset of $(HH_0({{\mathbf A}},{{\mathbf A}}g)^s)^*$ determined by the nondegenera\-cy condition, which is the vanishing condition for countably many determinant polynomials.
\end{proof}

\begin{Remark} \quad
\begin{enumerate}\itemsep=0pt
\item It is clear from the proof of Proposition~\ref{finde1} that it holds more generally, namely when~$g$ acts trivially on its generalized 1-eigenspace (this is the only consequence of semisimplicity of~$g$ used in the proof).
\item In particular, Proposition \ref{finde1} applies to the setting of \cite{L2}, namely $A=\mathcal{O}(X)$, where~$X$ is a conical symplectic singularity in the sense of Beauville. Indeed, in this case it is known that~$X$ has finitely many symplectic leaves \cite[Theorem~2.3]{Ka}. This property also holds for Higgs branches considered in~\cite{BPR} (by \cite[Section~2.3]{L3}), and is expected for Coulomb branches which are cones, so Proposition~\ref{finde1} applies to all such cases.
\end{enumerate}
\end{Remark}

\setcounter{Theorem}{13}
% Original source lines 506--545
\subsubsection{The general case}
Now let us generalize Proposition \ref{finde1} to the nonsemisimple case. For this purpose assume that every Poisson derivation of $A$ of nonpositive degree is inner. Then the Lie algebra $\g$ of the group
$G:=\operatorname{Aut}_0(A)$ of filtration-preserving Poisson automorphisms of $A$ commuting with $s$ has the form $\g=A_2$, where the operation is the Poisson bracket.

Let us also assume that the group~$G$ is reductive. Then it is easy to see that the action of $\g$ on $A$ uniquely lifts to any quantization ${\mathbf A}$, so we have ${\rm Der}({\mathbf A})^s=\g$. Thus we have a Lie algebra inclusion $\g\subset F_2{\mathbf A}\subset {\mathbf A}$.

Let $g\in G$. Write $g=g_0\exp(e)=\exp(e)g_0$ where $g_0\in G$ is semisimple and $e\in \g$ is nilpotent.

\begin{Proposition}\label{finde2} Under the above assumptions, the conclusion of Proposition~{\rm \ref{finde1}} holds for~$g$. Moreover,
\begin{gather*}
\dim HH_0({\mathbf A},{\mathbf A}g)^s\le \dim HH_0({\mathbf A},{\mathbf A}g_0)^s\le \dim HP_0(X^{g_0})^s<\infty.
\end{gather*}
\end{Proposition}

\begin{proof} Let $Z$ be the centralizer of $g_0$ in $G$, a reductive group containing $\exp(e)$. By the Jacobson--Morozov lemma, the Lie algebra $\mathfrak{z}$ of $Z$ contains an $\mathfrak{sl}_2$-triple
$e$, $h$, $f$. We have a~decomposition ${\mathbf A}=\oplus_{j\in {\mathbb Z}}{\mathbf A}_j$, where ${\mathbf A}_j$ is the $j$-eigenspace of~$h$. Since $g(e)=e$, we have $e{\mathbf b}-{\mathbf b}g(e)=[e,{\mathbf b}]$, so for any $T\in HH_0({\mathbf A},{\mathbf A}g)^*$ and ${\mathbf b}\in {\mathbf A}$ we have $T([e,{\mathbf b}])=0$. By representation theory of $\mathfrak{sl}_2$, this implies that $T({\mathbf a})=0$ for any ${\mathbf a}\in {\mathbf A}_j$ with $j>0$.

Now consider the 1-parameter subgroup $t^h$ of $G$ ($t\in {\mathbb C}^\times$), and define
$T_t({\mathbf a}):=T\big(t^{-h}{\mathbf a}t^h\big)$. If ${\mathbf a}\in {\mathbf A}_j$ then $T_t({\mathbf a})=t^{-j}T({\mathbf a})$. Thus, since $T$ vanishes on ${\mathbf A}_j$ for $j>0$, there exists a limit $T_0=\lim\limits_{t\to 0}T_t$, which coincides with $T$ on ${\mathbf A}_0$ and vanishes on ${\mathbf A}_j$ for $j\ne 0$.

We have $T({\mathbf a})=T_t({\mathbf a}_t)$, where ${\mathbf a}_t:=t^h{\mathbf a} t^{-h}$.
Recall that
\begin{gather*}
{\mathbf a}{\mathbf b}-{\mathbf b} g({\mathbf a})={\mathbf a}{\mathbf b}-{\mathbf b}g_0({\mathbf a})-{\mathbf b}\sum_{i\ge 1}\frac{g_0e^i}{i!}({\mathbf a})\in \operatorname{Ker}(T).
\end{gather*}
Hence
\begin{gather*}
t^h({\mathbf a}{\mathbf b}-{\mathbf b} g({\mathbf a}))t^{-h}={\mathbf a}_t{\mathbf b}_t-{\mathbf b}_t g_0({\mathbf a}_t)-{\mathbf b}_t\sum_{i\ge 1}\frac{t^{2i}e^i}{i!}({\mathbf a}_t)\in \operatorname{Ker}(T_t).
\end{gather*}
Thus, for any ${\mathbf a},{\mathbf b}\in {\mathbf A}$ the element ${\mathbf a}{\mathbf b}- {\mathbf b} g_0({\mathbf a})-{\mathbf b}\sum\limits_{i\ge 1}\frac{t^{2i}g_0e^i}{i!}({\mathbf a})$ belongs to $\operatorname{Ker}(T_t)$. Sending~$t$ to zero, we see that ${\mathbf a}{\mathbf b}-{\mathbf b} g_0({\mathbf a})\in \operatorname{Ker}(T_0)$. Thus, $T_0\in (HH_0({\mathbf A},{\mathbf A}g_0)^s)^*$, which is finite dimensional (with dimension dominated by $\dim HP_0(X^{g_0})^s$) by Proposition~\ref{finde1}.

It remains to show that the assignment $T\mapsto T_0$ is injective; this implies both statements of the theorem. To this end, assume that $T_0=0$ but $T\ne 0$, and let $m>0$ be the order of vanishing of $T_t$ at $t=0$. Let $T_0':=\lim\limits_{t\to 0}T_t/t^m$. Then $T_0'$ is a nonzero element of $HH_0({\mathbf A},{\mathbf A}g_0)^*$ which vanishes on ${\mathbf A}_0$. But if ${\mathbf a}\in {\mathbf A}_j$, $j\ne 0$, then ${\mathbf a}=\frac{1}{j}[h,{\mathbf a}]$, so $T_0'({\mathbf a})=0$. This is a~contradiction, which completes the proof of the proposition.
\end{proof}

\begin{Theorem}\label{finde3} Let $A=\mathcal{O}(X)$, where $X$ is a conical symplectic singularity
with Poisson bracket of degree $-2$ such that the group $G$ of dilation-invariant Poisson automorphisms of $X$ is reductive. Then for any quantization ${\mathbf A}$ of $A$ and any $g\in G$ the space $HH_0({\mathbf A},{\mathbf A}g)^s$ is finite dimensional, with dimension dominated by $\dim HP_0(X^{g_0})^s$, which is finite. Therefore, nondegenerate short star-products on~$A$ depend on finitely many parameters.
\end{Theorem}

\begin{proof}The theorem follows immediately from the above results and Lemmas \ref{inn}, \ref{finfix}.
\end{proof}
\begin{thebibliography}{99}
\bibitem[1]{BPR}
Beem C., Peelaers W., Rastelli L., Deformation quantization and superconformal
 symmetry in three dimensions, \href{https://doi.org/10.1007/s00220-017-2845-6}{\textit{Comm. Math. Phys.}} \textbf{354} (2017),
 345--392, \href{https://arxiv.org/abs/1601.05378}{arXiv:1601.05378}.

\bibitem[6]{Sie}
de~Siebenthal J., Sur les groupes de {L}ie compacts non connexes,
 \href{https://doi.org/10.1007/BF02564352}{\textit{Comment. Math. Helv.}} \textbf{31} (1956), 41--89.

\bibitem[11]{ES}
Etingof P., Schedler T., Poisson traces and {$D$}-modules on {P}oisson
 varieties, \href{https://doi.org/10.1007/s00039-010-0085-4}{\textit{Geom. Funct. Anal.}} \textbf{20} (2010), 958--987, \href{https://arxiv.org/abs/0908.3868}{arXiv:0908.3868}.

\bibitem[16]{HKLR}
Hitchin N.J., Karlhede A., Lindstr\"om U., Ro\v{c}ek M., Hyper-{K}\"{a}hler
 metrics and supersymmetry, \href{https://doi.org/10.1007/BF01214418}{\textit{Comm. Math. Phys.}} \textbf{108} (1987), 535--589.

\bibitem[17]{Ka}
Kaledin D., Symplectic singularities from the {P}oisson point of view,
 \href{https://doi.org/10.1515/CRELLE.2006.089}{\textit{J.~Reine Angew. Math.}} \textbf{600} (2006), 135--156,
 \href{https://arxiv.org/abs/math.AG/0310186}{arXiv:math.AG/0310186}.

\bibitem[19]{L2}
Losev I., Deformations of symplectic singularities and orbit method for
 semisimple {L}ie algebras, \href{https://arxiv.org/abs/1605.00592}{arXiv:1605.00592}.

\bibitem[20]{L3}
Losev I., Bernstein inequality and holonomic modules, \href{https://doi.org/10.1016/j.aim.2016.12.033}{\textit{Adv. Math.}}
 \textbf{308} (2017), 941--963, \href{https://arxiv.org/abs/1501.01260}{arXiv:1501.01260}.

\bibitem[23]{Lu}
Lusztig G., Character sheaves on disconnected groups.~{I}, \href{https://doi.org/10.1090/S1088-4165-03-00204-8}{\textit{Represent.
 Theory}} \textbf{7} (2003), 374--403, {E}rrata,
 \href{https://doi.org/10.1090/S1088-4165-04-00247-X}{\textit{Represent.
 Theory}} \textbf{8} (2004), 179, \href{https://arxiv.org/abs/math.RT/0305206}{arXiv:math.RT/0305206}.

\bibitem[25]{Na3}
Namikawa Y., Poisson deformations of affine symplectic varieties, {II},
 \href{https://doi.org/10.1215/0023608X-2010-012}{\textit{Kyoto~J. Math.}} \textbf{50} (2010), 727--752, \href{https://arxiv.org/abs/0902.2832}{arXiv:0902.2832}.

\bibitem[26]{Na1}
Namikawa Y., Fundamental groups of symplectic singularities, in Higher
 Dimensional Algebraic Geometry~-- in Honour of {P}rofessor {Y}ujiro
 {K}awamata's Sixtieth Birthday, \textit{Adv. Stud. Pure Math.}, Vol.~74,
 Editors K.~Oguiso, C.~Birkar, S.~Ishii, S.~Takayama, \href{https://doi.org/10.2969/aspm/07410321}{Math. Soc. Japan}, Tokyo,
 2017, 321--334, \href{https://arxiv.org/abs/1301.1008}{arXiv:1301.1008}.

\bibitem[30]{Ste}
Steinberg R., Conjugacy classes in algebraic groups, \textit{Lecture Notes in
 Math.}, Vol.~366, \href{https://doi.org/10.1007/BFb0067854}{Springer-Verlag}, Berlin~-- New York, 1974.

\end{thebibliography}
\end{document}
